    \subsection{Overview of DTR PCS \DTRPcs}

    Checks for compliant behavior of the generator under a grid scenario where there is a
    symmetric three-phase fault at one of the four transmission lines (at 1\%
    distance from the PDR).

    The grid model and its operational point is as in the following schematic:
    \begin{center}
        \includestandalone[width=0.8\textwidth]{circuit_\DTRPcs.tikz}
    \end{center}
    \begin{center}
        \small \textbf{Note: This schematic is only a reminder of the test setup on the TSO's
        side --- the Producer's side may vary, depending on the user-provided model.}
    \end{center}

    \noindent Important reminders:
    \begin{itemize}
        \item The time for fault clearance, $T_{clear}$, is a parameter that
        is specified in the DTR and depends on the voltage level in
        question:
        \begin{itemize}
            \item HTB3 (400 kV): $T_{clear}$ = 85 ms
            \item HTB2 (150 kV and 225 kV): $T_{clear}$ = 85 ms
            \item HTB1 (63 kV and 90 kV): $T_{clear}$ = 150 ms
        \end{itemize}
    \end{itemize}

    \begin{description}
        \item Initial conditions used at the PDR bus:
        \begin{itemize}
            \item $P = P_{max\_unite}$
            \item $Q = 0$
            \item $U = U_{dim}$
        \end{itemize}
    \end{description}

    \subsubsection{Simulation parameters}

    Solver and parameters used in the simulation:
    \begin{center}
        \begin{tabular}{lc}
            \toprule
           \textbf{Parameter} & \textbf{Value (default)} \\
            \midrule
            \BLOCK{for row in solverPCSI4ThreePhaseFaultTransientBolted}
            {{row[0]}}         & {{row[1]}}                         \\
            \BLOCK{endfor}
            \bottomrule
        \end{tabular}
    \end{center}

    \subsubsection{Simulation}
    As required by the DTR PCS \DTRPcs, the figures below show the
    following magnitudes:
    \begin{itemize}
        \item Voltage at connection point: modulus of the AC complex voltage at
        the PDR bus.
        \item Reactive power supplied at the connection point.
        \item Active power supplied at the connection point.
        \item Rotor speed.
        \item Magnitude controlled by the AVR.
    \end{itemize}

    \subsubsection{Simulation results}
    The blue line shows the calculated curve, if a reference curve has been entered it is
    shown in orange.

% For now we won't use floats for figures, to get more precise placement
    \noindent
    \begin{minipage}[t]{0.48\textwidth}
        \centering
        \includegraphics[width=\textwidth]{\Producer_fig_V_PCS_RTE-I4.ThreePhaseFault.TransientBolted}
        \begin{minipage}[t]{0.8\textwidth}
            \small Voltage magnitude measured at the PDR bus.
            The vertical dashed line marks the rise time $T_{85U}$.
        \end{minipage}
    \end{minipage}
    \hfill
        \begin{minipage}[t]{0.48\textwidth}
            \centering
            \includegraphics[width=\textwidth]{\Producer_fig_Ustator_PCS_RTE-I4.ThreePhaseFault.TransientBolted}
            \begin{minipage}[t]{0.8\textwidth}
                \small Stator voltage magnitude. The gray dotted line show
                the AVR setpoint.
            \end{minipage}
        \end{minipage}

    \vspace{0.5cm}

    \noindent
    \begin{minipage}[t]{0.48\textwidth}
        \centering
        \includegraphics[width=\textwidth]{\Producer_fig_P_PCS_RTE-I4.ThreePhaseFault.TransientBolted}
        \begin{minipage}[t]{0.8\textwidth}
            \small Real power output P, measured at the PDR bus. Green-dashed
            lines: $\pm 10\%$ margin around the final value of P. Red-dashed
            lines: same, for $\pm 5\%$ margin.
        \end{minipage}
    \end{minipage}
    \hfill
    \begin{minipage}[t]{0.48\textwidth}
        \centering
        \includegraphics[width=\textwidth]{\Producer_fig_Q_PCS_RTE-I4.ThreePhaseFault.TransientBolted}
        \begin{minipage}[t]{0.8\textwidth}
            \small Reactive power output Q, measured at the PDR bus.
        \end{minipage}
    \end{minipage}

    \vspace{0.5cm}

    \begin{minipage}[t]{0.48\textwidth}
        \centering
        \includegraphics[width=\textwidth]{\Producer_fig_W_PCS_RTE-I4.ThreePhaseFault.TransientBolted}
        \begin{minipage}[t]{0.8\textwidth}
            \small Rotor speed.
        \end{minipage}
    \end{minipage}
    \\[2\baselineskip]
    Go to  {{ linkPCSI4ThreePhaseFaultTransientBolted }}


    \subsubsection{Analysis of results}

    \noindent Analysis of results for PCS \DTRPcs. Values extracted
    from the simulation:

    \begin{center}
        \begin{tabular}{lc}
            \toprule
            \textbf{Parameter} & \multicolumn{1}{c}{\textbf{value}} \\
            \midrule
            \BLOCK{for row in rmPCSI4ThreePhaseFaultTransientBolted}
            {{row[0]}}         & {{row[1]}}                         \\
            \BLOCK{endfor}
            \bottomrule
        \end{tabular}
    \end{center}

    \noindent Key:
    \begin{description}
        \item[Settling time $T_{10P}$:] time at which the suplied real power
        P stays within the +/-10\% tube centered on the final value of P.
        \item[Settling time $T_{5P}$:] time at which the suplied real power
        P stays within the +/-5\% tube centered on the final value of P.
        \item[Rise time $T_{85U}$:] time at which the voltage at the PDR
        bus returns back above 0.85 pu, regardless of any possible
        overshooting / undershooting that may take place later on.
        \item[Critical clearing time $T_{cct}$:] the time threshold
        for the fault to clear, before the onset of rotor-angle stability.
    \end{description}


    \subsubsection{Compliance checks}

    Compliance checks required by PCS \DTRPcs:
    \begin{center}
        \begin{tabular}{lcc}
            \toprule
            \textbf{Check} & \multicolumn{1}{c}{\textbf{value}} & \multicolumn{1}{c}{\textbf{Compliant?}} \\
            \midrule
            \BLOCK{for row in cmPCSI4ThreePhaseFaultTransientBolted}
            {{row[0]}}     & {{row[1]}}                         & {{row[2]}}                              \\
            \BLOCK{endfor}
            \bottomrule
        \end{tabular}
    \end{center}
